\documentclass[11pt,a4paper]{article}
\usepackage[utf8]{inputenc}
\usepackage{amsmath,amssymb,amsthm}
\usepackage{algorithm}
\usepackage{algpseudocode}
\usepackage{booktabs}
\usepackage{hyperref}
\usepackage{geometry}
\geometry{margin=1in}

\title{Bidirectional Recurrent Neural Networks (BiRNN): Dual-Temporal Unrolling, Contextual Fusion, and Information Invariance}
\author{Team Scaibu \\ \small Email: \texttt{jaydeep.9664920749@gmail.com} \\ \small Architecture and Core Engine Team}
\date{\today}

\newtheorem{theorem}{Theorem}
\newtheorem{lemma}[theorem]{Lemma}
\newtheorem{proposition}[theorem]{Proposition}
\theoremstyle{definition}
\newtheorem{definition}{Definition}
\newtheorem{remark}{Remark}

\begin{document}

\maketitle

\begin{abstract}
Standard unidirectional recurrent architectures process temporal sequence inputs in a strictly causal, forward direction, restricting the hidden state at step $t$ to historical context $x_{1:t}$ and creating blind spots regarding subsequent future tokens $x_{t+1:T}$. Bidirectional Recurrent Neural Networks (BiRNNs) resolve this causal constraint by running two decoupled recurrent passes in opposite temporal directions: a forward network processing $1 \to T$ and a backward network processing $T \to 1$. Concatenating these dual representations yields an omnidirectional contextual encoding for offline sequence labeling. This document formalizes the bidirectional state-space model, proves contextual mutual information upper bounds, specifies algorithmic pseudocode, and walks through a verified numerical calculation.
\end{abstract}

\section{Notation and Mathematical Preliminaries}

Let $\mathcal{X} = (\mathbf{x}_1, \dots, \mathbf{x}_T)$ denote the complete observed input sequence of length $T$, where each token $\mathbf{x}_t \in \mathbb{R}^{d_x}$.

\begin{table}[h]
\centering
\caption{Mathematical Notation for Bidirectional Recurrent Networks}
\begin{tabular}{lll}
\toprule
\textbf{Symbol} & \textbf{Domain} & \textbf{Semantic Description} \\
\midrule
$\overrightarrow{\mathbf{h}}_t$ & $\mathbb{R}^{d_f}$ & Forward recurrent hidden state at time $t$ \\
$\overleftarrow{\mathbf{h}}_t$ & $\mathbb{R}^{d_b}$ & Backward recurrent hidden state at time $t$ \\
$\mathbf{h}_t$ & $\mathbb{R}^{d_f + d_b}$ & Fused bidirectional contextual representation: $[\overrightarrow{\mathbf{h}}_t \,\|\, \overleftarrow{\mathbf{h}}_t]$ \\
$\overrightarrow{W}_x, \overrightarrow{W}_h$ & Matrices & Forward input and recurrence parameter matrices \\
$\overleftarrow{W}_x, \overleftarrow{W}_h$ & Matrices & Backward input and recurrence parameter matrices \\
$W_y, \mathbf{b}_y$ & Matrix, Vector & Output emission projection parameters \\
\bottomrule
\end{tabular}
\end{table}

\section{Problem Definition and Core Concepts}

\subsection{3-Level Explanation}
\begin{enumerate}
    \item \textbf{Intuitive (High School level):} If you are reading a sentence with a missing word like ``The \underline{\hspace{1cm}} barked loudly at the mailman'', reading only forward gives you no clue what word fits. Reading both forward (seeing ``The'') and backward (seeing ``barked at the mailman'') instantly tells you the missing word is ``dog''.
    \item \textbf{Engineering Level:} BiRNN executes two separate recurrent state updates: forward states $\overrightarrow{\mathbf{h}}_t$ propagate from $t=1$ to $T$, while backward states $\overleftarrow{\mathbf{h}}_t$ propagate from $t=T$ down to $1$. At every time step $t$, the two hidden vectors are concatenated into $[\overrightarrow{\mathbf{h}}_t; \overleftarrow{\mathbf{h}}_t]$.
    \item \textbf{Mathematical Level:} Given sequence $\mathbf{x}_{1:T}$, the state recurrence is governed by the decoupled system:
    \begin{align}
    \overrightarrow{\mathbf{h}}_t &= \tanh(\overrightarrow{W}_h \overrightarrow{\mathbf{h}}_{t-1} + \overrightarrow{W}_x \mathbf{x}_t + \overrightarrow{\mathbf{b}}_h), \quad \overrightarrow{\mathbf{h}}_0 = \mathbf{0} \\
    \overleftarrow{\mathbf{h}}_t &= \tanh(\overleftarrow{W}_h \overleftarrow{\mathbf{h}}_{t+1} + \overleftarrow{W}_x \mathbf{x}_t + \overleftarrow{\mathbf{b}}_h), \quad \overleftarrow{\mathbf{h}}_{T+1} = \mathbf{0} \\
    \mathbf{h}_t &= [\overrightarrow{\mathbf{h}}_t \,\|\, \overleftarrow{\mathbf{h}}_t] \in \mathbb{R}^{d_f + d_b} \\
    \mathbf{y}_t &= W_y \mathbf{h}_t + \mathbf{b}_y
    \end{align}
\end{enumerate}

\subsection{5-Step Algorithmic Framework}
\begin{enumerate}
    \item \textbf{Sequence Verification:} Validate non-empty sequence dimension $T \ge 1$.
    \item \textbf{Forward Pass Execution:} Iterate $t = 1 \dots T$, updating and saving forward hidden states $\overrightarrow{\mathbf{h}}_t$.
    \item \textbf{Backward Pass Execution:} Iterate $t = T \dots 1$, updating and saving backward hidden states $\overleftarrow{\mathbf{h}}_t$.
    \item \textbf{Contextual Concatenation:} Form fused vectors $\mathbf{h}_t = [\overrightarrow{\mathbf{h}}_t; \overleftarrow{\mathbf{h}}_t]$ for each position $t \in \{1,\dots, T\}$.
    \item \textbf{Sequence-Level Emission:} Multiply concatenated features by emission weights $W_y$ to yield per-step predictions.
\end{enumerate}

\section{Algorithm Specification}

\begin{algorithm}[H]
\caption{Bidirectional Recurrent Neural Network (BiRNN)}
\label{alg:birnn}
\begin{algorithmic}[1]
\Require Sequence $\{\mathbf{x}_t\}_{t=1}^T$, forward parameters $(\overrightarrow{W}_x, \overrightarrow{W}_h, \overrightarrow{\mathbf{b}}_h)$, backward parameters $(\overleftarrow{W}_x, \overleftarrow{W}_h, \overleftarrow{\mathbf{b}}_h)$, emission parameters $(W_y, \mathbf{b}_y)$.
\Ensure Bidirectional representations $\{\mathbf{h}_t\}_{t=1}^T$ and predictions $\{\mathbf{y}_t\}_{t=1}^T$.
\State $\overrightarrow{\mathbf{h}}_0 \gets \mathbf{0}, \quad \overleftarrow{\mathbf{h}}_{T+1} \gets \mathbf{0}$
\For{$t = 1$ \textbf{to} $T$} \Comment{Forward Recurrence}
    \State $\overrightarrow{\mathbf{h}}_t \gets \tanh(\overrightarrow{W}_h \overrightarrow{\mathbf{h}}_{t-1} + \overrightarrow{W}_x \mathbf{x}_t + \overrightarrow{\mathbf{b}}_h)$
\EndFor
\For{$t = T$ \textbf{downto} $1$} \Comment{Backward Recurrence}
    \State $\overleftarrow{\mathbf{h}}_t \gets \tanh(\overleftarrow{W}_h \overleftarrow{\mathbf{h}}_{t+1} + \overleftarrow{W}_x \mathbf{x}_t + \overleftarrow{\mathbf{b}}_h)$
\EndFor
\For{$t = 1$ \textbf{to} $T$} \Comment{Fusion and Emission}
    \State $\mathbf{h}_t \gets [\overrightarrow{\mathbf{h}}_t \,\|\, \overleftarrow{\mathbf{h}}_t]$
    \State $\mathbf{y}_t \gets W_y \mathbf{h}_t + \mathbf{b}_y$
\EndFor
\State \Return $\{\mathbf{h}_t\}_{t=1}^T, \{\mathbf{y}_t\}_{t=1}^T$
\end{algorithmic}
\end{algorithm}

\section{Theoretical Guarantees and Contextual Coverage}

\begin{theorem}[Total Contextual Mutual Information Bound]
\label{thm:birnn_context}
Let $I(\mathbf{h}_t; \mathcal{X})$ denote the mutual information between the hidden representation at step $t$ and the complete sequence $\mathcal{X}$. Under the BiRNN parameterization:
\begin{equation}
I(\mathbf{h}_t; \mathcal{X}) = I(\overrightarrow{\mathbf{h}}_t; \mathbf{x}_{1:t}) + I(\overleftarrow{\mathbf{h}}_t; \mathbf{x}_{t:T}) - I(\overrightarrow{\mathbf{h}}_t; \overleftarrow{\mathbf{h}}_t \mid \mathcal{X})
\end{equation}
guaranteeing non-zero statistical dependence on both causal past ($\mathbf{x}_{1:t}$) and future context ($\mathbf{x}_{t:T}$).
\end{theorem}

\begin{proof}
By chain rule for mutual information, $I([\mathbf{A}; \mathbf{B}]; \mathcal{X}) = I(\mathbf{A}; \mathcal{X}) + I(\mathbf{B}; \mathcal{X} \mid \mathbf{A})$. Since $\overrightarrow{\mathbf{h}}_t$ is a deterministic function of $\mathbf{x}_{1:t}$ and $\overleftarrow{\mathbf{h}}_t$ is a deterministic function of $\mathbf{x}_{t:T}$, Markov property yields $I(\overrightarrow{\mathbf{h}}_t; \mathbf{x}_{t+1:T} \mid \mathbf{x}_{1:t}) = 0$, giving the stated decomposition.
\end{proof}

\section{Complexity Analysis}

\begin{itemize}
    \item \textbf{Time Complexity:} Exactly $2\times$ the FLOPs of a unidirectional RNN: $\mathcal{O}(2 T \cdot (d_h^2 + d_h d_x))$.
    \item \textbf{Space Complexity:} $\mathcal{O}(2 T \cdot d_h)$ memory to store forward and backward trajectories.
\end{itemize}

\section{Exactness and Error Analysis}

BiRNN representations require the entire sequence $\mathcal{X}$ to be available offline prior to starting the backward pass. Hence, BiRNN cannot be used for online streaming or autoregressive causal generation where future tokens are unobserved.

\section{Worked Numerical Example}

Let $T=2, d_x=1, d_f=1, d_b=1$. Parameters:
$\overrightarrow{W}_x = [1.0], \overrightarrow{W}_h = [0.5], \overrightarrow{b} = [0.0]$.
$\overleftarrow{W}_x = [1.0], \overleftarrow{W}_h = [0.5], \overleftarrow{b} = [0.0]$.
Input sequence: $x_1 = 1.0, x_2 = -1.0$.

\textbf{Forward Pass:}
\begin{itemize}
    \item $\overrightarrow{h}_1 = \tanh(0.5(0) + 1.0(1.0)) = \tanh(1.0) \approx 0.7616$.
    \item $\overrightarrow{h}_2 = \tanh(0.5(0.7616) + 1.0(-1.0)) = \tanh(0.3808 - 1.0) = \tanh(-0.6192) \approx -0.5506$.
\end{itemize}

\textbf{Backward Pass:}
\begin{itemize}
    \item $\overleftarrow{h}_2 = \tanh(0.5(0) + 1.0(-1.0)) = \tanh(-1.0) \approx -0.7616$.
    \item $\overleftarrow{h}_1 = \tanh(0.5(-0.7616) + 1.0(1.0)) = \tanh(-0.3808 + 1.0) = \tanh(0.6192) \approx 0.5506$.
\end{itemize}

\textbf{Concatenated Representations:}
\begin{itemize}
    \item $h_1 = [0.7616, 0.5506]$.
    \item $h_2 = [-0.5506, -0.7616]$.
\end{itemize}

\section{Numerical Considerations and Edge Cases}

\begin{itemize}
    \item \textbf{Padding Masking:} In batched variable-length sequences, the backward RNN must initialize at the true sequence length $L_i \le T_{\max}$ rather than padding tokens to prevent absorbing zero-padding noise into backward states.
\end{itemize}

\section{References}
\begin{enumerate}
    \item Schuster, M., & Paliwal, K. K. (1997). Bidirectional recurrent neural networks. \emph{IEEE Transactions on Signal Processing}, 45(11), 2673--2681.
    \item Graves, A., & Schmidhuber, J. (2005). Framewise phoneme classification with bidirectional LSTM and other neural network architectures. \emph{Neural Networks}, 18(5-6), 602--610.
\end{enumerate}

\end{document}
